\documentclass[11pt]{article}
\usepackage[T1]{fontenc}
\usepackage[utf8]{inputenc}
\usepackage{amsmath,amssymb,amsthm}
\usepackage[margin=2.5cm]{geometry}
\usepackage{booktabs}
\usepackage{graphicx}
\setlength{\parindent}{1.4em}
\setlength{\parskip}{0pt}

\newtheorem{theorem}{Theorem}
\newtheorem{definition}{Definition}
\newtheorem{corollary}{Corollary}
\newtheorem{remark}{Remark}
\newtheorem{proposition}{Proposition}

\title{\Large\textbf{Platform Information Laundering:}\\
How Low-Cost Aggregated Reactions Become High-Cost Validation Metrics}
\author{Hakvin Vosteen\\\small Universit\"at Bremen\\\small \texttt{vosteen@uni-bremen.de}}
\date{2026}

\begin{document}

\maketitle

\begin{abstract}
\noindent Classical signaling theory \cite{spence1973} assumes that observers process signals efficiently given their information sets. This assumption is plausible for professional observers (venture capitalists, banks, recruiters) whose economic incentives justify costly information processing. It fails for the dominant observer class on modern social platforms: retail users who react impulsively, at near-zero cost, and without verification. We argue that platforms create an information laundering mechanism by aggregating low-cost retail reactions into metrics (likes, follower counts, engagement scores) that are subsequently read by professional observers as if they were the output of high-cost professional verification. This pricing arbitrage between two distinct information-processing modes explains why signal-substance decoupling has accelerated in markets with strong platform presence. The mechanism does not require new behavioral assumptions on individual agents; it follows from the architectural transformation of unaggregated impulses into aggregated metrics. We develop a formal model, identify three observer classes, derive the laundering pipeline, and discuss empirical signatures and policy implications. Bibliometric evidence from the OpenAlex API (2015--2025) shows research output on platform-information pathologies growing $4.0\times$ faster than baseline economics, social science, and psychology output, consistent with the predicted structural acceleration. The framework provides a micro-foundation for the platform component of contemporary LARP \cite{vosteen2026larp} and connects to the macro pattern of regulatory waterbed effects \cite{vosteen2026waterbed}.
\end{abstract}

\section{Introduction}

The classical \cite{spence1973} model of signaling has two protagonists: the agent who produces costly signals and the observer who infers type from the signal. The model assumes the observer processes signals in a Bayesian-efficient manner, updating beliefs optimally given the available information. This assumption is reasonable when the observer has the technical capacity to interpret the signal and the economic incentive to invest the necessary processing effort. Professional financial intermediaries, hiring committees at top firms, and regulators with adequate funding approximately satisfy these conditions.

The assumption fails sharply on modern social platforms. The dominant observer on TikTok, Instagram, LinkedIn, and Twitter is the retail user who scrolls through content for entertainment, encounters a signal in a few seconds, and reacts with a single-tap operation (like, share, follow) before continuing to scroll. The retail user has neither the processing capacity of a venture capitalist nor the economic incentives that would justify investing in such capacity. Their reaction is, in any meaningful sense, an impulse.

The architectural innovation of platforms is to aggregate millions of such impulses into single metrics: a like count, a follower count, an engagement rate. These aggregated metrics are then visible to a different class of observer: professional decision-makers who use them as inputs to high-stakes resource allocation. Venture capitalists examine LinkedIn follower counts when evaluating founders. Publishers examine Twitter follower counts when evaluating authors. Recruiters examine engagement metrics when evaluating candidates. Universities examine social media presence when evaluating faculty applicants.

The platform thus operates as a transformation layer: it takes input from a class of observers who cannot process signals (retail users) and produces output that is read by a class of observers who could process signals (professionals) as if it were professional output. We call this process \emph{information laundering}, by direct analogy to financial laundering, in which money of dubious origin is processed through a network of small transactions until it acquires the appearance of legitimate funds. Here, signals of dubious substance are processed through a network of low-cost reactions until they acquire the appearance of validated quality.

The mechanism is not behavioral. We do not assume that retail users are stupid, or that professional observers are lazy. The mechanism is purely architectural: the aggregation step interposes between the retail observer and the professional observer, and the aggregation collapses the distinction between informed and uninformed reaction. After aggregation, a like from a domain expert and a like from a casual scroller are indistinguishable. The platform produces a metric that contains no information about who liked or why.

This paper formalizes the mechanism, identifies its empirical signatures, and connects it to the broader literature on signaling-mechanism failure. Section 2 introduces the three observer classes. Section 3 formalizes the laundering operator. Section 4 derives empirical implications. Section 5 connects to \cite{vosteen2026larp} on LARP at the micro level and \cite{vosteen2026waterbed} on regulatory waterbed effects at the macro level. Section 6 discusses policy implications. Section 7 acknowledges limitations.

\section{Three Observer Classes}

\begin{definition}[Observer Classes]
\label{def:classes}
We distinguish three observer classes in modern signaling markets:
\begin{itemize}
\item \textbf{Class A (Professional observer):} High processing capacity, strong economic incentives to invest in information acquisition. Examples: venture capitalists evaluating founders, banks evaluating borrowers, recruiters at top firms evaluating candidates.

\item \textbf{Class B (Retail observer):} Low processing capacity, weak or absent economic incentives. Examples: platform users who encounter content during entertainment-driven scrolling. Reactions are impulsive single-tap operations without verification.

\item \textbf{Class C (Aggregated metric):} Not an observer in the agentic sense, but a derived quantity produced by aggregating reactions from Class B. Examples: total likes, follower counts, engagement rates, virality scores.
\end{itemize}
\end{definition}

\begin{center}
\includegraphics[width=0.85\textwidth]{figures/observer_classes.pdf}
\end{center}

The central architectural fact about platforms is the following:

\begin{proposition}[Misreading of C as A]
\label{prop:misreading}
Class A observers systematically read Class C metrics as informative signals about underlying type, despite the fact that C is produced by Class B observers who do not perform type-inference.
\end{proposition}

The misreading is not necessarily irrational from the Class A observer's perspective. If the observer believes that high-quality content is more likely to attract attention, they may treat aggregated attention as a noisy but informative signal. The point is that this belief is approximately correct only when the underlying reactions are themselves informative. When the reactions are impulsive, the aggregate carries little information about quality, only about attention-capture capacity, which is a different and largely decoupled property.

\begin{remark}
A common defense of platform metrics is that ``the wisdom of crowds'' aggregates many imperfect signals into an accurate average. This defense requires that individual signals be approximately unbiased estimators of the underlying quantity. Social media reactions are not unbiased: they systematically over-weight emotional content, visual salience, in-group identification, and engagement-optimized format. The aggregate is biased in the direction of these features rather than toward underlying substance.
\end{remark}

\section{The Laundering Operator}

\begin{definition}[Laundering Operator]
\label{def:laundering}
A platform laundering operator $\mathcal{L}$ maps a tuple $(s, \mathbf{r})$ to an aggregated metric $M$, where $s$ is the agent's signal and $\mathbf{r} = (r_1, \ldots, r_n)$ are the impulsive reactions of $n$ Class B observers:
\[
\mathcal{L}: (s, \mathbf{r}) \mapsto M = \sum_{i=1}^{n} r_i.
\]
\end{definition}

The operator is a pure aggregation. It does not condition on the identity, expertise, or processing depth of the contributing reactors. Two key properties follow:

\begin{proposition}[Information Compression]
\label{prop:compression}
The map $\mathcal{L}$ is a strict information compressor: the Shannon entropy $H(M)$ is strictly less than the entropy $H(\mathbf{r})$ of the raw reaction vector. Information about the distribution of reactor types is destroyed by aggregation.
\end{proposition}

\begin{proof}[Sketch]
The map $\mathbf{r} \mapsto \sum r_i$ is many-to-one. Any two reaction vectors with the same sum produce the same $M$. The compression is therefore strict.
\end{proof}

\begin{proposition}[Substance Decoupling]
\label{prop:decoupling}
Let $\theta$ denote underlying agent type and $s$ the agent's platform signal. The conditional mutual information $I(M; \theta \mid s)$ is bounded above by the conditional mutual information $I(\mathbf{r}; \theta \mid s)$, with equality only when the aggregation preserves all information relevant to type inference, which is generically not the case.
\end{proposition}

Together, Propositions \ref{prop:compression} and \ref{prop:decoupling} establish that the aggregated metric $M$ contains strictly less information about $\theta$ than the underlying reaction vector $\mathbf{r}$. In practice, the gap is large because the relevant information for type inference (which reactors are expert? how carefully did they evaluate?) is the information that aggregation specifically destroys.

\begin{center}
\includegraphics[width=0.95\textwidth]{figures/laundering_pipeline.pdf}
\end{center}

\subsection{The Five-Step Pipeline}

The complete laundering pipeline operates in five steps:

\begin{enumerate}
\item \textbf{Signal production.} L-type agent posts content $s$ at near-zero marginal cost. The content need not be backed by underlying substance.

\item \textbf{Impulse reaction.} Class B observers react to $s$ based on superficial features (visual salience, emotional resonance, format-optimization), producing reactions $r_i$ without verification.

\item \textbf{Aggregation.} Platform sums reactions to produce $M = \sum r_i$.

\item \textbf{Misreading.} Class A observer reads $M$ as a substance signal, interpreting it as the output of informed type-inference rather than the aggregation of impulsive reactions.

\item \textbf{Resource extraction.} Class A observer allocates resources (funding, hiring, publishing, etc.) to the agent based on $M$. The agent extracts value without producing substance.
\end{enumerate}

The pipeline can run continuously and at large scale. There is no natural correction mechanism within the pipeline itself, because each step is locally rational: Class B observers react how they always react, the platform aggregates how it always aggregates, and Class A observers read the metric in the way that has historically been informative in non-platform contexts.

\section{Empirical Implications}

The framework yields testable predictions about when and where signal-substance decoupling should be most severe.

\subsection{Prediction 1: Cross-Platform Variation}

Platforms differ in the architecture of their aggregation. TikTok aggregates over a broader and more impulsive reactor population than LinkedIn, which in turn aggregates over a broader population than peer-reviewed publication systems. The laundering effect should scale with the breadth and impulsivity of the underlying reactor population.

Specifically: the L-type extraction success rate, measured as the rate at which platform-validated agents fail to deliver substance once funded, should be highest in TikTok-validated markets, intermediate in LinkedIn-validated markets, and lowest in peer-review-validated markets.

\subsection{Prediction 2: Decoupling Acceleration with Platform Penetration}

Pre-2010, professional resource allocators relied primarily on professional verification channels (peer review, references, direct evaluation). Post-2015, they increasingly rely on platform metrics as preliminary screens. The laundering effect should therefore have accelerated over the period 2010-2025, with measurable effects on the substance-validation gap.

This prediction is consistent with the rise in high-profile cases of platform-validated agents whose substance was discovered to be absent (vanity-metric founders, influencer-economy scandals, citation-farming academics).

\subsection{Prediction 3: Domain Sensitivity}

Domains in which professional observers can quickly verify substance (open-source software, formal mathematical proofs) should show weaker laundering effects than domains where verification is structurally delayed (biotech, defense, private equity). In low-verification domains, the time window during which platform metrics dominate professional verification is long enough for substantial resource extraction.

\subsection{Prediction 4: Architectural Sensitivity}

Platform design choices that reduce aggregation breadth (e.g., showing reactor identity, requiring substantive engagement to react, exposing reaction-time distributions) should reduce the laundering effect. The empirical test is to compare platforms before and after such architectural changes.

\subsection{Bibliometric Evidence}

If platform information laundering is a structural mechanism that has accelerated over the last decade, we expect academic research attention on platform-information pathologies to grow faster than baseline academic output. This is the revealed-preference logic developed in \cite{vosteen2026waterbed}: research agendas track problems the academic community perceives as growing, so disproportionate growth in a research area is indirect evidence that the underlying problem is itself growing.

We query the OpenAlex API (queried 2026-05-14) for the annual count of works matching eight platform-laundering-adjacent terms over 2015--2025: engagement-metrics reliability, viral content dynamics, influencer credibility, social media misinformation, platform algorithmic amplification, vanity metrics, parasocial trust, and opinion polarization on platforms. As a baseline we compute the same statistics for three broad fields (economics, social science, psychology) and use their mean growth as the comparison reference.

\begin{center}
\includegraphics[width=0.98\textwidth]{figures/bibliometric_evidence.pdf}
\end{center}

\medskip\noindent\textbf{Results.} The baseline growth rate of broad academic output over 2015--2025 is $3.3\times$. Seven of the eight laundering-adjacent topics grow at $4.7\times$ to $30.5\times$, exceeding the baseline rate by factors of $1.4\times$ to $9.1\times$. Mean excess growth is $4.0\times$ the baseline rate; the median is $3.8\times$. The strongest excess growth is in research on parasocial trust ($9.1\times$), engagement-metrics reliability ($5.8\times$), and platform algorithmic amplification ($4.9\times$). Spearman correlations of annual counts with year are at or near $\rho = 1.00$ for all eight series (all $p < 0.05$). Welch $t$-tests comparing pre-2020 to post-2020 means find significant differences ($p < 0.05$) for six of the eight series; the two non-significant cases (engagement-metrics reliability, vanity metrics) reflect either large variance in early years or small absolute counts.

\medskip\noindent\textbf{Interpretation.} The bibliometric signature is consistent with the existence of a structural platform-information mechanism that has become salient to the academic community over the last decade. We emphasize the limits of this evidence: bibliometric growth is consistent with multiple causal stories (genuine problem growth, fashion-driven attention shifts, post-COVID digital salience, post-2016 misinformation salience). The data do not isolate the specific Class-A misreading mechanism from these alternatives. They establish only that the topic cluster is growing structurally faster than the broader fields, which is consistent with the framework and inconsistent with the null hypothesis that platform information dynamics are a minor or stable component of contemporary information economics.

\medskip\noindent The bibliometric methodology applied here is the same one used in \cite{vosteen2026waterbed} for the CSRD-Waterbed analysis, where it provided indirect evidence for the Asymmetry Migration Principle. The consistency of the methodology across both papers, applied to two different domains (sustainability disclosure regulation and platform information dynamics), provides a unified empirical foundation for the broader Asymmetry-Migration framework.

\section{Connections to Adjacent Frameworks}

\subsection{Micro: LARP Theory}

\cite{vosteen2026larp} develops a theory of strategic LARP (life-action role-playing) at the micro level, distinguishing four LARP modes (Mimicry, Volume, Status, Inversion) keyed to observer cognitive hierarchy level $k$. The platform laundering mechanism complements this theory by explaining how the cognitive hierarchy distribution of observers became bimodal in the platform era.

Before platforms: most signal observation was performed by Class A observers operating at $k = 2$ or higher. LARP required either substance mimicry against naive observers (Mimicry-LARP) or sophisticated inversion against expert subcultures (Inversion-LARP).

After platforms: signal observation is performed by Class B observers at $k = 0$ (impulsive reaction), with the output read by Class A observers who treat the $k = 0$ aggregate as if it were $k = 2$ output. The effective observer cognitive level seen by the LARPer is $k = 0$, even when the resource-allocating observer would have operated at $k = 2$ in direct evaluation.

The laundering mechanism is therefore not merely additive to LARP theory. It systematically shifts the relevant observer cognitive level downward, making Volume-LARP (which exploits $k = 0$ aggregation) far more profitable than it would be without the laundering bridge.

\subsection{Macro: Regulatory Waterbed}

\cite{vosteen2026waterbed} documents the regulatory waterbed effect for mandatory disclosure regulation: when one channel of information asymmetry is regulated, the asymmetry migrates to adjacent channels. The platform laundering mechanism is a particular instance of this migration. When professional verification channels (peer review, audited disclosure, regulated reporting) compress information asymmetry, agents migrate to platform channels where laundering remains available. The macro pattern of asymmetry migration finds its concrete micro implementation in the platform laundering pipeline.

\section{Policy Implications}

The framework yields several policy implications.

\subsection{For Platform Design}

Aggregation breadth is a policy variable. Platforms that aggregate uniformly over all reactors maximize laundering. Platforms that condition the visibility of metrics on the verified expertise of contributing reactors reduce laundering. Specific design choices that reduce the laundering effect:

\begin{itemize}
\item Display reactor identity or credentialing alongside aggregated metrics.
\item Require minimum engagement time before reactions are counted in aggregate metrics.
\item Expose reaction-time distributions so professional observers can detect impulse-dominated metrics.
\item Provide expert-weighted alternative aggregates alongside raw counts.
\end{itemize}

None of these eliminate laundering. They reduce its effectiveness by reducing the information loss in the aggregation step.

\subsection{For Professional Observers}

The framework suggests a practical heuristic for Class A observers: when evaluating an agent based on platform metrics, discount the metrics by the laundering coefficient
\[
\lambda = 1 - \frac{I(M; \theta \mid s)}{I(\mathbf{r}; \theta \mid s)},
\]
which measures the fraction of type-relevant information destroyed by aggregation. In practice, $\lambda$ approaches one for platforms with broad and impulsive reactor populations, indicating that almost no type-information survives aggregation.

\subsection{For Regulators}

The platform laundering mechanism operates outside the scope of traditional disclosure regulation. SEC disclosure rules do not apply to LinkedIn metrics. FDA disclosure rules do not apply to TikTok claims. The information laundering occurs in a regulatory gap between platform regulation (which addresses content moderation but not metric integrity) and financial regulation (which addresses disclosure but not platform metrics).

Closing this gap would require explicit regulation of the metric-validation interface: rules about what platform metrics can be claimed as evidence of underlying substance in regulated contexts (e.g., investment prospectuses, employment claims, academic credentials).

\section{Limitations}

The model treats Class B reactions as uniformly impulsive. In practice, some platform users do perform thoughtful evaluation, particularly in expert communities. The model approximates the dominant pattern rather than the universal pattern.

The model treats the aggregation function as a simple sum. Real platforms use weighted aggregations, machine-learned ranking algorithms, and conditional visibility rules. The qualitative result (information compression) is robust to these complications, but the quantitative magnitude depends on aggregation details.

The model does not formally separate reactor heterogeneity from temporal dynamics. Some platforms have evolving reactor populations (older platforms accumulate expert users; newer platforms attract impulsive users). The cross-platform predictions in Section 4 may be confounded by these dynamics.

The model treats Class A observers as exogenously naive about the laundering mechanism. In reality, sophisticated investors are increasingly aware of platform metric unreliability and adjust their evaluation accordingly. The model describes the regime in which laundering succeeds, not the asymptotic regime in which professional observers fully discount platform metrics.

\section{Conclusion}

Classical signaling theory assumes that observers process signals efficiently. This assumption holds for professional observers in low-platform-penetration markets. It fails systematically in modern markets where platform architecture interposes an aggregation layer between low-cost retail reactors and high-cost professional decision-makers. The aggregation layer destroys most of the information relevant to type inference, but the resulting metric is read by professional observers as if it were the output of informed verification.

The mechanism is architectural, not behavioral. It does not require irrationality on the part of any individual observer. It follows from the simple fact that aggregation of impulses produces a metric that is misclassified as the output of verification. The misclassification is rational in legacy contexts (where aggregated reactions did track substance) and incorrect in platform contexts (where they do not).

The framework provides a micro-foundation for the platform-era acceleration of signal-substance decoupling, complements the LARP theory of strategic mimicry at the agent level, and connects to the macro pattern of regulatory waterbed effects. It also identifies a specific class of policy interventions (aggregation-architecture rules) that target the mechanism directly rather than its symptoms.

The deeper implication is that platforms are not neutral distribution layers for information. They are information-transformation layers that systematically transform low-cost impulses into high-cost-appearing validation metrics. The economics of modern signaling cannot be understood without explicit attention to this transformation.

\vspace{0.5cm}
\noindent\textbf{References}

\begin{thebibliography}{99}

\bibitem{akerlof1970} Akerlof, G. (1970). The market for lemons: Quality uncertainty and the market mechanism. \textit{Quarterly Journal of Economics} 84(3), 488--500.

\bibitem{benkler2006} Benkler, Y. (2006). \textit{The Wealth of Networks: How Social Production Transforms Markets and Freedom}. Yale University Press.

\bibitem{camerer2004} Camerer, C.F., Ho, T.-H., Chong, J.-K. (2004). A cognitive hierarchy model of games. \textit{Quarterly Journal of Economics} 119(3), 861--898.

\bibitem{goodhart1975} Goodhart, C.A.E. (1975). Problems of monetary management: The U.K. experience. \textit{Papers in Monetary Economics}, Reserve Bank of Australia.

\bibitem{hayek1945} Hayek, F.A. (1945). The use of knowledge in society. \textit{American Economic Review} 35(4), 519--530.

\bibitem{shannon1948} Shannon, C.E. (1948). A mathematical theory of communication. \textit{Bell System Technical Journal} 27(3), 379--423.

\bibitem{spence1973} Spence, M. (1973). Job market signaling. \textit{Quarterly Journal of Economics} 87(3), 355--374.

\bibitem{strathern1997} Strathern, M. (1997). Improving ratings: Audit in the British university system. \textit{European Review} 5(3), 305--321.

\bibitem{vosteen2026larp} Vosteen, H. (2026a). A theory of strategic LARP: Cognitive hierarchy, endogenous observation regimes, and type selection in signaling markets. SSRN Working Paper.

\bibitem{vosteen2026waterbed} Vosteen, H. (2026b). The regulatory waterbed: Why mandatory disclosure moves information asymmetry instead of removing it. SSRN Working Paper.

\end{thebibliography}

\end{document}
